\section{Bounded axes and specialization of effectivity}
\label{b:bounded-axes}

The width estimates still allow torus widths to approach zero, and their
upper counts involve the cost as well as the physical degree. This section
removes both difficulties in the range needed for the first optimization
step. Its key obstruction is a degree-zero section of bounded positive
cost, which would contradict concentration of the entire base maximal
ideal. To construct such a section from a hypothetical obstruction, we
must retain effectivity, not only numerical divisor classes.

Pass to a subsequence on which each width is either bounded above or tends
to infinity. Call the former axes \emph{cheap}, and let \(J\) be the set of
cheap torus indices. The cheap residual axes are precisely the \(p_f\) axes
from \(K_f\); there are \(r=|J|+p_f\) cheap axes in total.
Put \(F=Y_f\); this is \(S_f\) up to the small
\(\Q\)-factorialization in the last step of the tower. We use \(H_f\) for
its pullback to \(F\) as well. Its positive strict coefficients are bounded
below by a fixed positive constant, including after the final normalization
\(s_f=1\). All constants in this section are uniform along the retained
sequence; the finite covers and open restrictions introduced below may
depend on a subsequence.

For each \(j\in J\), the opposite probes constructed in
\eqref{eq:C5} and after \eqref{eq:C8} give
\[
 \bar h^{m_j/n}y_j^+,\ \bar h^{m_j/n}y_j^-\in R,
 \qquad \operatorname{wt}(y_j^\pm)=\pm d_j\beta_j.
\]
Here \(m_j\in n\Z_{\geq0}\) and the positive integers \(d_j\) are uniformly
bounded. Their boundedness follows from the upper bounds on the cheap
widths. The functions \(y_j^\pm\) are angular functions and have physical
degree zero.

Here is the estimate we will prove. For a homogeneous section of
\(R\), its physical degree is its \(D\)-degree \(m\), and its cost
\(b\) is its \(\tilde v\)-value. On an actual angular lift
\(F_0\in V_m\), the same cost is \(b=m+v(F_0)\).

\begin{proposition}[The bounded axes estimates]\label{b:bounded-axes-estimates}
All widths \(W_j\) are bounded below by a uniform positive constant.
There are cutoffs \(S_0\to\infty\) and a uniform constant \(C_4\) such
that
\begin{equation}\label{eq:D3}
\begin{gathered}
 b=\tilde v(\text{weight})\leq C_4m
       \qquad(0<b\leq S_0),\\
 \dim\bigl(\text{initial spaces at degree }m
                  \text{ through cost }S_0\bigr)
       \leq C_4(1+m)^r,
 \qquad r=|J|+p_f.
\end{gathered}
\end{equation}
The independent probes after \eqref{eq:C8} give matching box lower bounds,
up to uniform factors, for all the widths. There are \(r\) algebraically independent probes
of bounded degree and cost: the positive cheap torus probes and the
\(p_f\) residual probes from \(K_f\).
\end{proposition}

The proof is completed at the end of the section. Its main step is to
bound the cheap torus charges in terms of physical degree, as long as
both grow slowly compared with the collapsing scales. A failure will
produce a regular eigenfunction of physical degree zero and bounded
positive cost. We first reduce its regularity to a finite problem on
the terminal field \(K_f\).

\begin{lemma}[Descent of slow pure sections]\label{b:pure-descent}
Suppose, along the sequence, that
\[
 g\in K_1^*,\qquad \operatorname{div}_{S_1}(g)+MH_1\geq0,
 \qquad M\in\Q_{\geq0},\qquad Ms_{\exp}\longrightarrow0.
\]
Then, for all sufficiently late members,
\[
 g\in K_f^*,\qquad \operatorname{div}_{F}(g)+CMH_f\geq0
\]
for a uniform constant \(C\).
\end{lemma}

\begin{proof}
Maintain the inequality
\(\operatorname{div}(g)\geq-CM\Lambda_i\) on all divisorial valuations
of \(K_i\). It holds initially by effectiveness and pullback of the given
rational Cartier divisor. At a collapsing level, the pivot on the proper normal
arithmetic generic fibre has \(\Lambda_i\)-value \(s_i\). Every other strict
prime of that fibre comes from a strict prime on \(S_i\) with
\(H_i\)-coefficient zero. The effective trace boundary of \(Q_i\), when
present, gives \(\Lambda_i\leq\lambda_i\) there. Since
\(\lambda_i\lesssim s_i\) and \(Ms_i\to0\), every possible pole order of
\(g\) on this fibre is an integer bounded in absolute value by a quantity
tending to zero. Thus \(g\) has no poles on the fibre. Properness, normality
and the contraction property imply \(g\in K_{i+1}\).

For a divisorial valuation downstairs, choose a lift attaining the minimum
for \(C_i\) in \eqref{eq:B1}. Use
\(\Lambda_i\lesssim\lambda_{i+1}C_i\) and
\(\lambda_{i+1}\pi_\#C_i=\Lambda_{i+1}\). Dividing by the restriction
multiplicity gives the same maintained inequality at level \(i+1\).
There are at most dimension-many steps, so the resulting constant is
uniform. At the last level \(\lambda_f=O(s_{\exp})\), with the usual zero
convention when the tower has no collapsing step. Off \(\Supp H_f\), the
same integral-order argument excludes poles. On \(\Supp H_f\), the strict
coefficient gap and the effective trace of \(Q_f\) give
\(\Lambda_f\lesssim H_f\). This proves the claimed divisor inequality.
\end{proof}

Applying \eqref{eq:A9} and Lemma~\ref{b:pure-descent} to the products of
opposite probes gives
\[
 g_j:=y_j^+y_j^-\in K_f^*,\qquad
 \operatorname{div}_{F}(g_j)+CH_f\geq0.
\]
Let \(\Sigma\) be the union of \(\Supp H_f\) and the supports of the zero
divisors of the \(g_j\). The degrees of its components, and the degrees of
the zero and pole divisors of every \(g_j\), are bounded. Indeed the last
displayed inequality bounds the pole degrees; principality bounds the zero
degrees by the same quantity. The degree of \(H_f\) is bounded by its
moduli reserve and the effective Calabi--Yau trace, as in \eqref{eq:C3};
its strict coefficient gap bounds the degree of its support. Divisor data
are omitted when \(\dim F=0\).

For an invariant prime divisor \(E\) of \(\Spec R\), write
\(E|_{K_f}=cP\), with \(P\) normalized when the restriction is nontrivial,
and put \(\ell_j=E(y_j^+)\). At the trivial restriction put \(c=0\) and
interpret all restricted orders as zero. The last step of the tower and
regularity of the two probes give
\begin{equation}\label{b:finite-row-bounds}
 cC_f(P)\leq1-h_E/2\quad(c>0),\qquad
 -m_jh_E\leq\ell_j\leq cP(g_j)+m_jh_E.
\end{equation}
Consequently the rows \((c,h_E,(\ell_j)_{j\in J})\) belong to a fixed
finite set. To see this, \(c\) is integral and bounded because
\(C_f(P)\) has a uniform positive lower bound; also
\(h_E\in\frac1n\Z\cap[0,1]\) by \eqref{eq:A10}--\eqref{eq:A11}.
Applying \eqref{eq:C1} to the zero and pole divisors bounds
\(c|P(g_j)|\), using \eqref{b:finite-row-bounds}; finally each \(\ell_j\)
is integral. If \(h_E=0\), and either \(c=0\) or the centre of \(P\) is
not contained in \(\Sigma\), then all \(\ell_j=0\).

\subsection{A finite marked regularity problem}
Consider an occurring weight supported on the cheap torus axes,
\[
 m\chi+\sum_{j\in J}\alpha_j\beta_j,\qquad
 \left(m+\sum_{j\in J}|\alpha_j|\right)s_{\exp}\longrightarrow0.
\]
Raise a nonzero section of this weight to a bounded positive common power
clearing \(n\) and the \(d_j\). It becomes
\[
 \bar h^{m'/n}\prod_{j\in J}(y_j^+)^{a_j}g,
 \qquad n\mid m',\quad a_j\in\Z,\quad g\in K_1^*.
\]
Here \(m'\) is that common multiple of \(m\), and \(a_j\) is the same
multiple of \(\alpha_j/d_j\). Multiplying this expression by the opposite
regular probe to the power \(a_j\) when \(a_j>0\), and by the positive
regular probe to the power \(-a_j\) when \(a_j<0\), produces a pure section.
Its degree is \(m'+\sum|a_j|m_j\) and its angular coefficient is
\(g\prod_{a_j>0}g_j^{a_j}\). Equation~\eqref{eq:A9} and
Lemma~\ref{b:pure-descent} imply that \(g\in K_f\) and that its poles on
\(F\) lie in \(\Sigma\). Regularity of the original eigenfunction is
exactly the family of inequalities
\begin{equation}\label{eq:D1}
 cP(g)+m'h_E+a\cdot\ell\geq0.
\end{equation}
These tests suffice because a pole divisor of a rational eigenfunction is
invariant, and \(R\) is normal.

We shall exclude an obstruction sequence with
\(\sum|a_j|/m'\to\infty\), counting \(m'=0\), \(a\ne0\) as infinite.
Suppose such a sequence exists and first assume \(\dim F>0\).
Our aim is a single rational polyhedral test for the possible degrees
and charges. The finitely many row types do not yet give such a test:
the projected prime valuations \(P\) and the existence of an effective
representative for the required pole divisor still depend on the
member. We next place the relevant primes on bounded marked models.

\begin{lemma}[A bounded guard model]\label{b:guard-model}
After passing to a subsequence, there are bounded projective
\(\Q\)-factorial varieties of Fano type \(F'\to F\) and labelled prime
divisors on \(F'\), of bounded number and degree, with the following
properties. Every nontrivial projection in \eqref{eq:D1} for which
\(h_E>0\), or whose centre is contained in \(\Sigma\), is one of the
marked prime valuations. Every component of \(\Sigma\) has marked strict
transform, and every exceptional prime of \(F'\to F\) is marked. The
pullback of \(H_f\) has bounded degree on \(F'\).
\end{lemma}

\begin{proof}
The effective Calabi--Yau datum \(C_f\) has a uniform positive discrepancy
bound, a boundary dominating a fixed positive multiple of \(H_f\), and a
big nef b-part with a fixed positive \(\mathbf H_f\)-reserve. Choose a
fixed rational \(C_0\) such that
\[
 \mathbf T_j=\operatorname{div}(g_j)+C_0\mathbf H_f
\]
is the b-Cartier closure of an effective divisor on \(F\). For sufficiently
small fixed rational \(\varepsilon>0\), define guard data by
\[
 C^{\mathrm g}=C_f-\varepsilon\sum_{j\in J}\mathbf T_j.
\]
Subtract the same terms from the moduli presentation. Equation~\eqref{eq:C1}
bounds \(\mathbf T_j(P)\) by a uniform multiple of \(C_f(P)\), so this
retains a uniform positive discrepancy bound. Principal terms do not affect
nefness, and the reserve absorbs \(\varepsilon C_0|J|\mathbf H_f\).
Thus the new moduli remain big and nef, with a positive
\(\mathbf H_f\)-reserve, and the trace boundary remains effective.

There is a fixed \(e_0>0\) such that all the prime valuations specified in
the statement satisfy \(C^{\mathrm g}(P)\leq1-e_0\). For a projection
with \(h_E>0\), this follows from \eqref{b:finite-row-bounds}, since
\(h_E\geq1/n\) and \(c\geq1\). If \(h_E=0\), Equation~\eqref{eq:B2}
gives \(P(H_f)=0\). A centre contained in \(\Sigma\) therefore lies
generically outside \(\Supp H_f\) and inside a zero divisor of some
\(g_j\); its integral order \(P(g_j)\) is at least one. Since
\(C_f(P)\leq1\), subtracting \(\varepsilon\mathbf T_j\) gives the
required gap. A strict prime of \(F\) in \(\Supp H_f\) has a boundary
coefficient bounded below by domination and the strict coefficient gap.
A remaining strict prime of \(\Sigma\) has the same positive zero order
just used. This also proves the assertion for all components of \(\Sigma\).

Apply Lemmas~\ref{b:ordinary-approximation} and
\ref{b:controlled-extraction}. The monomial calculation in the proof of
Lemma~\ref{b:scale-bad-primes} gives only finitely many exceptional
prime valuations with \(C^{\mathrm g}(P)\leq1-e_0<1\), including
equality at the upper endpoint: their positive integral normal weights
satisfy a weighted-sum bound below one, and each coefficient in that
sum is positive.
Ordinary approximation with a tiny ample reserve does not increase their
discrepancies. The extraction construction gives a projective
\(\Q\)-factorial model \(F'\to F\) extracting exactly these primes.
The guard trace is effective on \(F'\), retains its uniform positive
log discrepancy bound, and has big nef moduli. Ordinary approximation
therefore gives an effective big boundary with rationally trivial adjoint
and a uniform positive discrepancy bound. By
\cite[Corollary~1.4]{BirkarBAB}, the varieties \(F'\) are bounded; they
are also of Fano type by the ample-reserve version of the construction.

Mark every strict prime of \(F'\) whose guard discrepancy is at most
\(1-e_0\). Each has guard boundary coefficient at least \(e_0\).
The trace degree of the nef b-part is nonnegative by the perturbation
argument in \eqref{eq:C2}. Thus the degree of the effective guard boundary
is at most the degree of \(-K_{F'}\), which is uniformly bounded for the
bounded very ample polarizations by the cone-theorem calculation preceding
\eqref{eq:C1}. This bounds both the number and the degrees of the marks.
The surviving \(\mathbf H_f\)-reserve gives the asserted bound on the
degree of the pullback of \(H_f\) in the same way.
\end{proof}

Label the marks along a subsequence, preserving which are exceptional and
which are strict over components of \(\Sigma\). Stabilize the finite sets
of rows in \eqref{eq:D1} projecting to each mark and the rows with trivial
projection. Introduce a rational signed pole allowance \(r_P\) at each
mark: a function \(g\) is allowed when \(P(g)\ge-r_P\), so a
negative allowance requires a zero. Consider the conditions
\begin{equation}\label{eq:D2}
\begin{gathered}
 m'\geq0,\qquad
 cr_P\leq m'h_E+a\cdot\ell
       \quad\text{for every row projecting to a mark }P,\\
 0\leq m'h_E+a\cdot\ell
       \quad\text{for every row with }c=0,\\
 r_P\leq0
       \quad\text{at marks strict over primes of }F\text{ outside }\Sigma,\\
 \sum_{P\ \mathrm{marked}}r_P[P]\in
       \overline{\operatorname{Eff}}(F').
\end{gathered}
\end{equation}
An actual candidate satisfies these conditions with \(r_P=-P(g)\):
\(\operatorname{div}_{F'}g+\sum r_PP\) is effective. Conversely, suppose
a rational solution has \(\sum r_PP\) rational-linearly equivalent to an
effective divisor. Choose a positive integer \(q\) and a function
\(f\in K_f^*\) with
\[
 \operatorname{div}_{F'}f+q\sum r_PP\ge0,
 \qquad qm'\in n\Z,\qquad qa\in\Z^J.
\]
The integer \(q\) clears the divisor coefficients and the degree and
charge denominators. For a marked row, effectivity and \eqref{eq:D2}
give
\[
 cP(f)+q(m'h_E+a\cdot\ell)
 \ge q\bigl(m'h_E+a\cdot\ell-cr_P\bigr)\ge0.
\]
The rows with trivial projection also remain valid after multiplication
by \(q\). The poles of \(f\) downstairs on \(F\) are confined to
\(\Sigma\): the sign condition excludes other strict marked primes,
and effectivity excludes all unmarked strict primes. Any remaining
nontrivial projection has \(h_E=0\), \(\ell=0\), and centre not
contained in \(\Sigma\), by Lemma~\ref{b:guard-model}. The function
is regular at the generic point of that centre, so its order is
nonnegative. Thus all tests \eqref{eq:D1} hold, and
\(\bar h^{qm'/n}\prod_j(y_j^+)^{qa_j}f\) is an actual regular
eigenfunction.

\subsection{A marked family with a Fano type generic fibre}
We now construct a common cone test for \eqref{eq:D2}. We need two
implications. Every actual candidate must satisfy the cone inequalities
of one geometric generic marked fibre, even when its coefficients
\(r_P\) are unbounded. Conversely, one fixed rational solution of
those inequalities must give an effective representative, and hence
an actual eigenfunction, on all sufficiently late retained members.
The first implication will use sweeping curves; the second will
spread a line bundle with a nonzero section.

Both arguments require the geometric generic marked fibre to be
\(\Q\)-factorial and of Fano type. We establish these properties by
retaining an ordinary klt Calabi--Yau boundary with a fixed denominator
that dominates a big marked divisor. The member-dependent generalized
guard data do not yet give this finite set of transferable divisor
coefficients and adjoint relations. The next lemma constructs it.

\begin{lemma}[A bounded auxiliary Calabi--Yau boundary]
\label{b:bounded-guard-boundary}
The models of Lemma~\ref{b:guard-model} admit effective integral divisors
\(P'\) of bounded degree, supported on the marks, such that \(|P'|\) is
birational. They also admit klt boundaries \(\Theta^+\), with bounded
degrees and coefficients in a fixed finite rational set, satisfying
\[
 \Theta^+\geq e'P',\qquad N(K_{F'}+\Theta^+)\sim0
\]
for a fixed positive rational \(e'\) and a bounded positive integer \(N\).
The marks, \(P'\), and \(\Theta^+\) together form a bounded family of
labelled divisor data.
\end{lemma}

\begin{proof}
Use the bounded \(k\) at the final level in \eqref{eq:B3} and set
\(P'=\lfloor kH_f\rfloor\) on \(F'\), using total pullback before
rounding. This is effective integral of bounded degree. Its support lies
in the marks, since the strict support of \(H_f\) and all exceptional
primes of \(F'\to F\) are marked. The functions in \eqref{eq:B3} still
belong to its complete system: their integral orders and positivity of
pullback allow the indicated rounding. Thus this system remains
birational, and \(P'\) is big.

Choose a small fixed rational \(e'>0\) so that
\[
 \Delta_0=e'\sum_{P\ \mathrm{marked}}P+e'P'
\]
is at most the guard boundary. Such a uniform choice is possible because
the mark coefficients have a fixed lower bound and the coefficients of
\(P'\) are bounded by its degree. By ordinary approximation of the big
nef b-part, obtain an effective uniformly positive-lc boundary \(\Theta\)
with
\[
 K_{F'}+\Theta\sim_{\Q}0,\qquad \Theta\geq\Delta_0,
 \qquad \Theta-\Delta_0\ \text{big}.
\]
Here the added effective representative of the big b-part has big trace:
sections on a descent remain sections after pushforward. It supplies the
last bigness assertion.

Set \(D'=-(K_{F'}+\Delta_0)\), which is big. Since \(F'\) is a
\(\Q\)-factorial variety of Fano type, it is a Mori dream space by
\cite[Corollary~1.3.1]{BCHM}.
Run a \(D'\)-MMP to a \(\Q\)-factorial model \(Y\) with nef and big
transform \(D_Y\). The program extracts no divisors. The rationally
trivial log adjoint \(K+\Theta\) is crepant through the program, so
\((Y,\Theta_Y)\), and hence \((Y,(\Delta_0)_Y)\), retain the uniform
positive discrepancy bound. Moreover
\(D_Y=-(K_Y+(\Delta_0)_Y)\) is nef and big. Thus \(Y\) is bounded by
\cite[Theorem~1.1]{BirkarBAB}. It is of Fano type: writing
\(D_Y\sim_{\Q}A+E\) with \(A\) ample and \(E\geq0\), a sufficiently
small positive multiple of \(E\) can be added to \((\Delta_0)_Y\) while
keeping klt, and the resulting negative log adjoint is ample.

The coefficients of \((\Delta_0)_Y\) have a fixed common denominator.
Lemma~\ref{b:class-lattices}, applied to the bounded \(\Q\)-factorial
varieties \(Y\), gives a bounded integer \(q>0\) with \(qD_Y\) Cartier.
Apply \cite[Theorem~1.1 and Remark~1.2]{FujinoBPF} to the klt pair
\((Y,(\Delta_0)_Y)\), with \(L=qD_Y\) and \(a=1\). The divisor
\[
 L-(K_Y+(\Delta_0)_Y)=(q+1)D_Y
\]
is nef and big; for a klt pair there are no additional lc-centre
conditions in log bigness. The theorem gives a bounded integer \(N>1\)
with \(ND_Y\) Cartier and free. Increase the free multiple if necessary.
For a general member
\(G_Y=\operatorname{div}_Y\psi+ND_Y\), Bertini on a log resolution
shows that
\((Y,(\Delta_0)_Y+G_Y/N)\) is klt.

Use the same rational canonical form on all models and define
\[
 \Theta^+=-K_{F'}+\frac1N\operatorname{div}_{F'}\psi.
\]
On a common resolution with maps \(p\) to \(F'\) and \(q_Y\) to \(Y\),
the MMP negativity comparison reads
\(p^*D'=q_Y^*D_Y+E'\), with \(E'\geq0\) exceptional over \(Y\).
Consequently
\[
 p^*(\Theta^+-\Delta_0)=q_Y^*(G_Y/N)+E'\geq0.
\]
Thus \(\Theta^+\geq\Delta_0\). Its log adjoint is crepant to that of
the klt pair just constructed on \(Y\), so it too is klt. Moreover
\(N(K_{F'}+\Theta^+)=\operatorname{div}_{F'}\psi\).
Its coefficients belong to \(\frac1N\Z\cap[0,1)\), and its degree is
the bounded degree of \(-K_{F'}\). In bounded projective embeddings,
the supports are parametrized by finitely many fixed-dimension,
fixed-degree Chow spaces \cite[Definition~3.21]{KollarChow2017}. Their
universal cycles may be taken over normal strata. The finite coefficient
set bounds the possible multiplicities, and the degree bound bounds the
number of prime components. Taking finite products and labelling their
components therefore bounds these data together with the marks and
\(P'\).
\end{proof}

Parametrize the labelled data of Lemma~\ref{b:bounded-guard-boundary} in
finite-type projective families over \(\C\), and pass to a subsequence
with fixed coefficient types and fixed \(N\). Choose a minimal closed
parameter subset containing infinitely many selected indices, then an
irreducible component \(T\) containing infinitely many of them. Every
proper closed subset of \(T\) contains only finitely many selected indices:
otherwise it would contradict minimality. Hence every dense open
restriction retains all but finitely many of those indices. Repetitions
of a parameter value cause no difficulty; in that case a minimal subset
may be a point. Denote the resulting marked family by
\(X^0\to T\), with fibres \(X^0_t\), marked primes, integral divisor
\(P'_t\), and boundary \(\Theta^+_t\).

This retention property survives the finite field extensions used below.
After normalization and restriction to a dense open, a finite extension of
\(\C(T)\) gives an irreducible finite surjective cover \(T'\to T\).
Every proper closed subset of \(T'\) has proper closed image, since finite
maps preserve dimension and both bases have the same dimension. Therefore
later dense open restrictions of \(T'\) exclude lifts only over a proper
closed subset of \(T\). Choose lifts of the remaining selected points.
The same reasoning applies after any finite composite of these extensions.

\begin{lemma}[The geometric generic marked model]
\label{b:generic-marked-fano}
After the finite covers and dense open restrictions just described, the
geometric generic marked fibre \(X^0_{\bar\eta}\) is
\(\Q\)-factorial and of Fano type. Its marked divisor \(P'_{\bar\eta}\)
is big, and its marked boundary \(\Theta^+_{\bar\eta}\) is klt with
\(N(K+\Theta^+_{\bar\eta})\sim0\).
\end{lemma}

\begin{proof}
We first transport the integral divisor-class relations to the generic
fibre. We then transport the crepant snc calculation for
\(\Theta^+\) and the section map of a Cartier multiple of \(P'\);
these establish klt singularities and bigness, respectively.

Spread a projective smooth resolution of the geometric generic fibre,
after a finite extension, and restrict the base so that it resolves every
fibre under consideration. Track all exceptional prime divisors and all
strict transforms of the finitely many marked primes. This uses the
resolution and constructibility argument in
Lemma~\ref{b:class-lattices}: keep a birational isomorphism on dense opens,
spread the complementary components, and fix their fibre and image
dimensions. Shrink also so that the downstairs fibres are geometrically
normal and integral and the prime marking families have integral
geometric fibres. Write this simultaneous resolution as
\(b_t:W_t\to X^0_t\).

On every selected fibre, vanishing for a log Fano boundary
\cite[Theorem~3.2]{FujinoFundamental} and rationality of klt
singularities \cite[Theorem~0.1]{FujinoRational} give
\(H^1(\OO)=H^2(\OO)=0\) both downstairs and upstairs. Semicontinuity and
density of the selected points give these vanishings throughout an open
base. On its complex fibres, the exponential sequence and GAGA
\cite[Section~12, Theorems~1--3; Section~20, Proposition~18]{SerreGAGA}
identify \(\Pic\) with integral \(H^2\). On a connected dense stratum, topological
constructibility gives locally constant integral group data
\[
 H^2(X^0_t,\Z)\ \xrightarrow{\ b_t^*\ }\ H^2(W_t,\Z)
 \ \longrightarrow\ H^2(W_t,\Z)/
      \langle[\text{exceptional primes}]\rangle.
\]
The last group is \(\Cl(X^0_t)\), and the image of the first group is
\(\Pic(X^0_t)\). The exceptional and marked upstairs divisor classes,
and the class of the relative canonical bundle of the smooth resolution,
are locally constant as well. These are integral diagrams, so their
cokernels, including torsion, are controlled; trivial monodromy is not
required.

On the selected fibres, the image of \(\Pic\) has finite index in
\(\Cl\), and \(N[K+\Theta^+]=0\) in \(\Cl\). The same assertions
therefore hold on this entire complex stratum. Here the canonical class
downstairs is the pushforward of the canonical class upstairs. They hold
also on the geometric generic fibre. For completeness, descend the finite
family, markings, resolution and open conditions to a countable
algebraically closed subfield of \(\C\), and choose a complex point
generic over that subfield. The geometric generic model is identified,
up to an isomorphism of algebraically closed fields preserving these
finite data, with this complex fibre: the two algebraically closed
extension fields have the same infinite transcendence degree over the
embedded defining function field. Thus the integral class conclusions
apply there. Finite index of \(\Pic\) in \(\Cl\) gives
\(\Q\)-factoriality, and the second conclusion gives the asserted
principal multiple of the log adjoint.

It remains to verify singularities and bigness; neither follows from the
integral class diagram alone. Take a log resolution of the geometric
generic marked pair, and a relative rational volume form and a rational
function witnessing \(N(K+\Theta^+)\sim0\). Spread all these data after
a further finite extension. On the smooth resolving family, retain the
divisor identities for the form and the witnessing function, the strict
transforms and exceptional components, and the snc pullback computation.
The coefficients and orders in these finite divisor formulas remain
unchanged on every fibre after shrinking. This can be arranged by spreading
each prime component of the formulas and retaining its geometric
integrality and its image dimension; restricting the formulas on the
smooth models then preserves the divisor orders. Pushing down preserves
the marked principal relation. The snc crepant boundary coefficients are now
constant, and one selected fibre is klt. They are therefore all less than
one on the geometric generic resolution, proving that
\(\Theta^+_{\bar\eta}\) is klt.

Choose a positive Cartier multiple \(tP'_{\bar\eta}\), and spread its
line bundle with a rational section whose divisor is exactly this marked
divisor. Retain the divisor identity by the same resolution argument.
After further shrinking, formation of its space of sections commutes with
base change. The dimension of the image of the associated rational map
is generically constant, by constructibility. On each selected fibre,
the system \(|tP'_t|\) contains products from the birational system
\(|P'_t|\), so its image has dimension \(\dim F'\). The generic system
has that same dimension. Hence \(P'_{\bar\eta}\) is big.

Finally \(\Theta^+_{\bar\eta}\geq e'P'_{\bar\eta}\) is big. To see
Fano type explicitly, write
\(\Theta^+_{\bar\eta}\sim_{\Q}A+E\) with \(A\) ample and \(E\geq0\).
For sufficiently small rational \(t>0\), the effective boundary
\((1-t)\Theta^+_{\bar\eta}+tE\) is klt, and its negative log adjoint
is rationally equivalent to \(tA\). This is a log Fano boundary.
\end{proof}

\subsection{Facet curves and one effective witness}

The geometric generic marked fibre now has the Fano type and rational
factoriality needed for its effective cone. We use this one cone in two
ways: finitely many sweeping curve families test every candidate on a
retained fibre, and one fixed rational solution supplies an effective
section that spreads to those fibres. These two implications suffice for
the degree-zero obstruction.

\begin{lemma}[Specializing a finite cone test]\label{b:marked-cone-test}
In \eqref{eq:D2}, replace the last condition by
\[
 \sum_{P\ \mathrm{marked}}r_P[P_{\bar\eta}]
       \in\overline{\operatorname{Eff}}(X^0_{\bar\eta}).
\]
This is a finite rational polyhedral test. After the allowed covers and
open restrictions, every actual candidate satisfies this test. Conversely,
each fixed rational solution has, after one clearing multiple fixed
for that solution, regular eigenfunction realizations on all retained
sufficiently late members of the sequence.
\end{lemma}

\begin{proof}
By Lemma~\ref{b:generic-marked-fano} and
\cite[Corollary~1.3.1]{BCHM}, \(X^0_{\bar\eta}\) is a Mori dream space.
Its pseudo-effective cone in numerical divisor space is rational
polyhedral. We first show that each facet is tested by a family of curves
sweeping a dense open of \(X^0_{\bar\eta}\).

Choose a rational divisor class \(D\) in the relative interior of that
facet. The class is pseudo-effective. Its divisor MMP ends in a minimal
model \(Y_D\) on which the transform \(D_Y\) is nef and semiample.
As \(D\) lies on the boundary of the pseudo-effective cone, \(D_Y\) is
not big. A sufficiently divisible multiple defines a contraction
\(\varphi:Y_D\to Z_D\) with positive-dimensional general fibres.
The birational program has no extraction. There is consequently an open
\(U\subset Y_D\), with complement of codimension at least two, isomorphic
to an open of \(X^0_{\bar\eta}\). On a common resolution the comparison
has the form
\[
 p^*D=q^*D_Y+E_D,\qquad E_D\geq0,
\]
where the correction is exceptional over \(Y_D\). Enlarge the excluded
codimension-two subset to contain its image and all inverse-map
indeterminacy. Curves in \(U\) then lift to the original variety and
avoid this correction, including the correction from any contracted
fixed divisorial part of \(D\).

A general fibre of \(\varphi\) is geometrically integral. Each component
of \(Y_D\setminus U\) either does not dominate \(Z_D\), and hence misses
a general fibre, or has codimension at least two in that fibre. If the
relative dimension is one, the latter case is impossible: a bad component
has dimension at most \(\dim Y_D-2<\dim Z_D\). Thus the entire general
curve fibre lies in \(U\). In larger relative dimension \(t\), intersect
a general fibre with \(t-1\) sufficiently general very ample hyperplanes.
Bertini gives an integral curve, and the codimension estimate makes it
miss the bad locus entirely.

These curves vary in a projective flat family over an integral parameter
space after restriction to a dense open of the fibre-and-hyperplane
parameter space. Its evaluation map dominates \(Y_D\). Indeed the
incidence obtained by choosing a point and hyperplanes through it is
irreducible and dominates \(Y_D\); restriction to the dense open of flat
integral curves avoiding the bad locus preserves dominance. The lifted
curves therefore sweep a dense open of \(X^0_{\bar\eta}\). Their
intersections with divisor classes are constant in this family. The
corresponding linear functional is nonnegative on every effective class,
because a member can be chosen not contained in the divisor's support;
it is positive on an ample class. Its value on \(D\) is zero by the
pullback comparison and contraction. A nonzero supporting functional
vanishing at a relative interior point of a facet defines exactly that
facet inequality. This constructs the required test for each facet.

There are finitely many facets. Spread the finitely many curve families,
their parameter spaces and evaluation maps after a finite extension of
the base function field. Spread also fixed Cartier multiples of all
marked divisors. After shrinking, the curves remain projective and flat
with integral geometric fibres, their intersection vectors with these
markings remain fixed, and their evaluation maps dominate dense opens of
every selected fibre. To justify the last assertion, the generic
evaluation image contains a dense open. Its image is constructible, so
after spreading and shrinking one retains an open in that image whose
intersection with each fibre is dense; fibre-dimension constructibility
excludes the exceptional parameter values. Equivalently one can track
the evaluation incidence and its generic dominance before the shrink.

Let \(\sum r_PP\) be the class furnished by an actual candidate on any
retained fibre. It has an effective rational Cartier representative.
For each of our curve families, density provides a member not contained
in that representative's support. Its intersection with the effective
representative is nonnegative, and is the fixed linear combination of
the marked intersection vector with coefficients \(r_P\). Thus every
generic facet inequality holds. This argument makes no boundedness
assumption on the coefficient vector of the candidate.

For the converse, fix one rational solution of the generic test.
On the generic Mori dream space, a rational class in the effective cone
has an effective rational representative. Numerical and rational linear
classes agree for rational Cartier divisors here by
Lemma~\ref{b:class-lattices}. Hence for some positive integer \(q\) there
is a rational function \(f\) with
\[
 \operatorname{div}(f)+q\sum r_PP_{\bar\eta}\geq0.
\]
Replace \((q,f)\) by \((kq,f^k)\) for a sufficiently divisible
positive integer \(k\), and retain the notation \((q,f)\). This
preserves the effective divisor inequality while making the marked
divisor Cartier and all degree and charge exponents integral, with
the degree divisible by \(n\). Spread the line bundle
\(\OO(q\sum r_PP_{\bar\eta})\), its rational trivialization, and the
nonzero regular section represented by \(f\), after a further finite
extension. The divisor identity can also be retained on the simultaneous
resolution. On a dense open of the parameter, the section is regular and
not identically zero on any fibre. Thus its rational-function
representatives have exactly the permitted pole bound on each selected
fibre; possible additional zeros are harmless. Regularity as a section
of this specified line bundle excludes unmarked poles.

The minimal-closure retention property ensures that these open
restrictions keep infinitely many selected indices, and eventually all
retained ones. Transfer the witness through the marked identifications
with the corresponding \(F'\). The converse argument following
\eqref{eq:D2}, including its automatic unmarked tests, gives the asserted
regular eigenfunctions. Only this one rational solution, its clearing
multiple and its effective section were spread. No simultaneous family
of the algebras \(R\) is needed, and no equality between the effective
cones of different fibres is asserted.
\end{proof}

\begin{proposition}[The slow-charge bound]\label{b:slow-charge-bound}
For sequences of actual candidates above satisfying
\( (m'+\sum|a_j|)s_{\exp}\to0\), the ratio
\(\sum|a_j|/m'\) cannot tend to infinity, with the convention already
specified at degree zero.
\end{proposition}

\begin{proof}
Suppose first that \(\dim F>0\). Use the stabilized rows of
\eqref{eq:D2} and the generic cone test of Lemma~\ref{b:marked-cone-test}.
They define a rational polyhedral cone in the variables \((m',a,(r_P))\).
Its projection \(\mathcal P\) to \((m',a)\) is again a rational
polyhedral cone and is therefore closed. For example, express the
original cone by finitely many rational generators and project those
generators. This proof of closedness imposes no boundedness condition
on the auxiliary coefficients \(r_P\).

Actual candidates lie in \(\mathcal P\). Divide an alleged obstruction
sequence by \(\sum|a_j|\). After a subsequence, their charge coordinates
converge to \(a_\infty\) of \(\ell^1\)-norm one, whereas their degree
coordinates tend to zero. Closedness gives
\((0,a_\infty)\in\mathcal P\). The rational face
\(\mathcal P\cap\{m'=0\}\) therefore contains a rational point with
nonzero charge: one can use its rational generators, or density of
rational points in its relative interior. A nonempty fibre of a rational
linear inequality system over a rational point has a rational point, so
choose rational auxiliary coordinates giving a fixed rational solution
of \eqref{eq:D2} with \(m'=0\) and \(a\ne0\).

By Lemma~\ref{b:marked-cone-test}, one fixed clearing multiple and one
spread effective witness produce actual nonzero regular eigenfunctions
on all sufficiently late retained rings \(R\), of degree zero and fixed
nonzero charge coordinates. If \(\dim F=0\), the same conclusion follows
directly: pure descent makes \(g\) constant, all restrictions to \(K_f\)
are trivial, and the stabilized trivial-projection inequalities in
\eqref{eq:D2} themselves give the rational polyhedral cone. There are no
markings or effectivity variables in this case.

The resulting eigenfunctions have positive \(\tilde v\)-cost, since a
nonzero charge is nonconstant and the grading is positive. These costs
are uniformly bounded: only cheap axes occur, their widths are bounded
above, the charge coordinates and the \(d_j\) are fixed or bounded, and
\(\tilde v\in Q_T\). Each eigenfunction lifts, in physical degree zero,
to a function of \(\mathcal A_0=k[Z]\) with the same positive value.
The lift therefore belongs to \(\mathfrak m_z\). Its bounded value
contradicts \eqref{eq:A7}, which asserts
\(\tilde v(\mathfrak m_z)\to\infty\) uniformly over all nonzero
functions in this ideal. This rules out the obstruction sequence.
\end{proof}

\subsection{The cost gap and the section count}

\begin{proof}[Proof of Proposition~\ref{b:bounded-axes-estimates}]
The residual widths \(1/s_i\) are already at least one. If some torus
width \(L_j\) tended to zero on a subsequence, choose integers
\(k\to\infty\) sufficiently slowly that \(kL_j\) remains bounded and
\(ks_{\exp}\to0\). Apply \eqref{eq:C5} to \(k\beta_j\), using a fixed
scalar norm bound at least one. It gives actual opposite weights of
bounded physical degree and charges \(\pm d_0k\beta_j\), with
\(d_0\) positive and bounded. These weights are supported on a cheap
axis, satisfy the slow condition, and have unbounded charge-to-degree
ratio. Raising to the earlier bounded common clearing power does not
change that conclusion. This contradicts
Proposition~\ref{b:slow-charge-bound}. A further-subsequence argument
therefore gives a uniform positive lower bound for every \(L_j\).

Choose \(S_0\to\infty\) so slowly that
\[
 S_0s_{\exp}\longrightarrow0,\qquad
 S_0/W_j\longrightarrow0\quad\text{for each width tending to infinity}.
\]
For a weight \(m\chi+\sum\alpha_j\beta_j\) of positive cost at most
\(S_0\), the weight comparisons preceding \eqref{eq:C7} give
\[
 m\lesssim S_0,\qquad
 \sum_j|\alpha_j|L_j\lesssim S_0.
\]
An expensive torus coordinate must vanish, since it is integral and its
width eventually exceeds every fixed multiple of \(S_0\). The lower
width bound then gives
\(m+\sum|\alpha_j|\lesssim S_0\), so these weights satisfy the slow
condition. Proposition~\ref{b:slow-charge-bound} now implies a uniform
\(\sum|\alpha_j|\lesssim m\) for all such weights: failure would give
an obstruction sequence on late members. The bounded clearing factors
and the bounded positive \(d_j\) preserve this equivalence.

Since the remaining widths are bounded above,
\[
 b=m+\tilde v\Bigl(\sum_{j\in J}\alpha_j\beta_j\Bigr)
 \leq m+\Bigl\|\sum_{j\in J}\alpha_j\beta_j\Bigr\|
 \lesssim m.
\]
Thus every initial at fixed degree \(m\) through \(S_0\) already occurs
through \(\min\{S_0,C_4m\}\). Apply \eqref{eq:C8} with this cutoff.
The factors for expensive widths are uniformly bounded, while each cheap
factor is at most a fixed multiple of \(1+m\). There are exactly
\(|J|+p_f=r\) cheap axes. This proves both estimates in \eqref{eq:D3},
including the constant initial at degree zero.

Finally the probe construction after \eqref{eq:C8} had box denominators
\(\max\{1,W_j\}\). The positive lower width bound makes these uniformly
comparable to \(W_j\). Hence those lower bounds now match all widths.
The residual probes have independent initials along the field tower, and
the positive torus probes add independent Laurent charges over \(K_1\).
Their combined initials are algebraically independent. The stated
\(r\) probes at the cheap axes have bounded degree and cost.
\end{proof}

Through the angular identification of Section~\ref{a:angular},
\eqref{eq:D3} controls the actual systems \(V_m\). At each degree
\(m\in n\Z_{\ge0}\), no function has cost in
\(C_4m<b(F)\le S_0\), and the quotient of \(V_m\) by its subspace
of cost greater than \(S_0\), including zero, has dimension at most
\(C_4(1+m)^r\). The \(r\) cheap probes have actual lifts in these
systems with bounded degree and cost and independent initials. These
are the finite-cutoff estimates used when varying the valuation on
\(k(X)\).
